% pack: with-previous
% TAB-27: generated by harness/analysis/paper_outputs.py; do not edit (every cell's source is in TAB-27.cells.json)
\begin{tabularx}{\linewidth}{l>{\hsize=0.536\hsize\raggedleft\hyphenpenalty=10000\exhyphenpenalty=10000\arraybackslash}X>{\hsize=0.807\hsize\raggedleft\hyphenpenalty=10000\exhyphenpenalty=10000\arraybackslash}X>{\hsize=0.923\hsize\raggedleft\hyphenpenalty=10000\exhyphenpenalty=10000\arraybackslash}X>{\hsize=1.552\hsize\raggedright\hyphenpenalty=10000\exhyphenpenalty=10000\arraybackslash}X>{\hsize=0.811\hsize\raggedleft\hyphenpenalty=10000\exhyphenpenalty=10000\arraybackslash}X>{\hsize=1.371\hsize\raggedright\hyphenpenalty=10000\exhyphenpenalty=10000\arraybackslash}X}
\toprule
Test & Locked $p$ (one-sided) & Inferiority $p$, one minus the locked $p$ (post hoc) & Holm threshold, inferiority direction (post hoc) & Inferiority $p$ below its Holm threshold (post hoc) & Stricter interval (post hoc) & Stricter interval wholly below the margin (post hoc) \\
\midrule
Test~1 & 0.99996 & \ensuremath{\leq}~0.00004 & 0.005\phantom{00} & yes & \mbox{[\ensuremath{-}47.50, \ensuremath{-}20.63]} & yes \\
Test~2 & 0.99993 & \ensuremath{\leq}~0.00007 & 0.00625 & yes & \mbox{[\ensuremath{-}52.23, \ensuremath{-}23.55]} & yes \\
Test~3 & 0.98913 & 0.01087 & 0.0125\phantom{0} & yes & \mbox{[\ensuremath{-}39.28, \ensuremath{-}8.25]} & no \\
Test~4 & 0.99987 & \ensuremath{\leq}~0.00013 & 0.00833 & yes & \mbox{[\ensuremath{-}57.50, \ensuremath{-}26.25]} & yes \\
Test~5 & \ensuremath{\leq}~0.00004 & 0.99996 & 0.025\phantom{00} & --- & --- & --- \\
\bottomrule
\end{tabularx}
% note: In the reading column, ``\ensuremath{>}~margin'' stands for ``by more than the margin'': ``worse by \ensuremath{>}~margin'' is worse by more than the margin, as the other tables of this supplement write it.
% note: Post hoc sensitivity, not a test: the correction was fixed before this check was computed, but ``chosen after the main results were open; the pre-registered unadjusted three-way reading stays the result'' (free analysis 5's own words; its ``three-way reading'' is the reading from the pre-registered decision rule, applied to the unadjusted interval).
% note: Holm in the inferiority direction (null hypothesis: $d$ at or above the margin of \ensuremath{-}10 points) over all five pre-registered tests at one-sided alpha 0.025; each inferiority $p$ is one minus the locked one-sided $p$, from the same bootstrap (same differences and seeds, $B$~=~10,000). Holm order, smallest inferiority $p$ first: Test~1, Test~2, Test~4, Test~3, Test~5.
% note: The stricter interval is the two-sided 99\% interval from the same bootstrap draws; it is shown beside Holm, not as part of it.
% note: Test~3 (post hoc): its inferiority $p$, 0.01087, is below its Holm threshold of 0.0125, but its 99\% interval \mbox{[\ensuremath{-}39.28, \ensuremath{-}8.25]} does not lie wholly below the margin of \ensuremath{-}10 points.
% note: Test~3 under the original marking rule (the sensitivity check Deviation~2 fixed in Lock~3, before any main-run data): $d$~=~+0.30 points, 95\% interval \mbox{[\ensuremath{-}13.58, 14.95]}; reading: inconclusive.
% note: Test~5's confirmatory result is its Holm-adjusted result in the pre-registered family of one-sided tests. Opus-low is shown no worse than Sonnet-medium by more than 10 points (Holm-adjusted): yes ($p$~\ensuremath{\leq}~0.00004, below its Holm threshold of 0.005). Sonnet ran at medium effort, one level below its default, in Test~3 and Test~5. Here Test~5 is only the fifth member of the inferiority-direction Holm family, not an inferiority claim; its full report is in ⟦TAB-02⟧.
% note: The error-rate ratio is the challenger's error rate over the comparator's, given beside $d$ as in ⟦TAB-02⟧ (pre-specified descriptive, reported beside the test, not part of its reading); the main analysis labels Test~1 and Test~2 as near ceiling (Amendment~1), where the error-rate ratio travels with $d$. Exact: the comparator's error-rate bound, in \%, is the upper end of its exact (Clopper--Pearson) 95\% interval.
% note: \ensuremath{\leq}~marks a $p$ at the smallest value the bootstrap can return (no resample reached the margin).
